\section{一混合Hesse三次式の完全resultant}
この節は形式比較から独立した三次式の幾何である。
\[
 F_{c,\kappa,\alpha}=(x^3+y^3+z^3)/3+cxyz+
 \kappa(\alpha uvx+u^3+v^3),\qquad d=1+c^3
\]
と置く。gradient resultantは$\operatorname{Res}(x^2,y^2,z^2,u^2,v^2)=1$で正規化する。
\begin{theorem}
全パラメータで多項式恒等式
\[
 \operatorname{Res}(F_x,F_y,F_z,F_u,F_v)
 =3^8\kappa^{32}d^3(\kappa\alpha^3+9)^3(\kappa\alpha^3+9d)^9
\]
が成立する。
$c\kappa d\ne0$なら特異条件は$(\kappa\alpha^3+9)(\kappa\alpha^3+9d)=0$であり、
前者は三つ、後者は九つの通常二重点を与える。
$c=0,\kappa\ne0$では二因子が一致して重複度十二となり、三特異点の局所型は
$p^2+q^2+y^3+z^3$である。
\end{theorem}
\begin{proof}
$\kappa d\ne0$の開集合でまず計算する。
$x=0$上では$F_y=F_z=0$から$y=z=0$、$F_u=F_v=0$から$u=v=0$なので射影的特異点はない。
$x=1$の勾配式は
\[
 y^2+cz=z^2+cy=0,\quad u^2+\alpha v/3=v^2+\alpha u/3=0,\quad
 1+cyz+\kappa\alpha uv=0.
\]
最初の対は$(0,0)$または$(-c\omega,-c\omega^2)$、$\omega^3=1$、
次の対は$(0,0)$または$(-\alpha\eta/3,-\alpha\eta^2/3)$、$\eta^3=1$である。
従って最後の式は二種類の因子を与え、点数はそれぞれ三、九となる。
四変数Hessianは二ブロック
\[
 \begin{pmatrix}2y&c\\c&2z\end{pmatrix},\qquad
 \kappa\begin{pmatrix}6u&\alpha\\\alpha&6v\end{pmatrix}
\]
に分かれ、その行列式はそれぞれ全体として$-3c^2\kappa^2\alpha^2$、
$9c^2\kappa^2\alpha^2$である。よって全点が通常二重点である。

定数と重複度をnormで決める。
$K[r,s]/(r^2+ts,s^2+tr)$の基底$(1,r,s,rs)$で$M_t=$乗法$\,rs$の行列は
\[
 M_t=\begin{pmatrix}0&0&0&0\\0&t^2&0&0\\0&0&t^2&0\\1&0&0&t^2\end{pmatrix}.
\]
従って上の四二次式による十六次元商上で$F_x$のnormは
\[
 \det(I_{16}+cM_c\otimes I_4+\kappa\alpha I_4\otimes M_{\alpha/3})
 =d^3(1+\kappa\alpha^3/9)^3(d+\kappa\alpha^3/9)^9.
\]
無限遠の先頭四式のresultantは$(3\kappa)^{16}$であり、
Poisson積公式\cite{Resultant}で次数二の$F_x$を加えるとその平方が掛かる。
分母$9^{12}$を整理して$3^8$を得る。
稠密開集合の計算は全パラメータの多項式恒等式へ延長する。
$c=0$では三つの$(u,v)$点でそのHessianは非退化で、$y,z$は独立三乗項のままなので表示の局所型となる。
$\kappa=0$なら頂点$\PP^1$を持つcone、$d=0$なら$u=v=0$上に三つの特異点が残ることも式から分かる。
\end{proof}

\begin{corollary}
$c=h^2/3,d=1+h^6/27,\kappa=d$とした候補式は
\[
 D/3+h^2xyz/3+d(\alpha uvx+u^3+v^3)
\]
であり、そのresultantは
$3^8d^{44}(d\alpha^3+9)^3(\alpha^3+9)^9$である。
$d\ne0$なら特異条件は
$(\alpha^3+9)((27+h^6)\alpha^3+243)=0$となる。
形式的には$b(h)^3=d,b(0)=1$により$u,v$を$b^{-1}$倍すれば
$D/3+h^2xyz/3+u^3+v^3+\alpha b(h)uvx$に移る。
この座標恒等式だけから有限$h$での圏同値や比較核の代数化は従わない。
\end{corollary}

\paragraph{有限係数候補の強制される規格化.}
$c=1/3$では$d=28/27$である。
この狭い形の候補で二つの悪い$\alpha^3$値を$-9,-243/28$へ一致させるには、
一般式の比が$d$であることから$\kappa=28/27$が強制される。
これは候補式のパラメータ照合であり、有限$h$での圏比較の存在証明ではない。
